\PassOptionsToPackage{quiet}{xeCJK}
\documentclass[answer]{THUExam}
\usepackage{HTNotes-math}
\usepackage{pifont}
\usepackage{ulem}

\usepackage{tikz}
\usetikzlibrary{positioning, calc, shapes.geometric, shapes, shapes.multipart, arrows.meta, arrows, decorations.markings, external, trees}

% =================================================
%       PDF信息
% =================================================
\title{高等数学A (二)}
\author{}
\Subject{作业}
\Keywords{高等数学}

% =================================================
%       试卷头信息
% =================================================
\Year{2025--2026}
\Semester{春季}
\Course{高等数学A}
\Time{}
\Type{A}

\begin{document}
\vspace{1cm}
\makehead

\vspace{1cm}

% =================================================
%       选择题：空间曲线、二次曲面与柱面旋转面
% =================================================
\makepart{选择题}{每小题~4~分，共~16~分}

\begin{problem}
在空间直角坐标系中，方程组
\[
\begin{cases}
x^2+y^2+z^2=4,\\
z=1
\end{cases}
\]
表示的图形是\pickout{B}
\options
{一个点}
{一条空间曲线}
{一个平面}
{一个球面}
\end{problem}
\begin{note}
空间曲线的一般方程可由两个空间曲面方程联立给出。此处 $z=1$ 是平面，$x^2+y^2+z^2=4$ 是球面，二者相交得到一个圆（属于空间曲线）。
\end{note}
\bigskip

\begin{problem}
方程 $\dfrac{x^2}{4}+\dfrac{y^2}{9}-\dfrac{z^2}{16}=1$ 所表示的二次曲面是\pickout{B}
\options
{椭球面}
{单叶双曲面}
{双叶双曲面}
{椭圆抛物面}
\end{problem}
\begin{note}
二次曲面的一般分类中，等号左边为两正一负的标准形式对应单叶双曲面；若为一正两负则为双叶双曲面；全正则对应椭球面。
\end{note}
\bigskip

\begin{problem}
xOy 平面上的椭圆 $\dfrac{x^2}{4}+\dfrac{y^2}{9}=1$ 绕 $x$ 轴旋转一周所得的旋转曲面方程为\pickout{B}
\options
{$\dfrac{x^2}{4}+\dfrac{y^2}{9}=1$}
{$\dfrac{x^2}{4}+\dfrac{y^2+z^2}{9}=1$}
{$\dfrac{x^2+y^2}{4}+\dfrac{z^2}{9}=1$}
{$\dfrac{x^2}{4}+\dfrac{y^2}{9}+\dfrac{z^2}{9}=1$}
\end{problem}
\begin{note}
曲线绕 $x$ 轴旋转时，将原方程中的 $y$ 替换为 $\pm\sqrt{y^2+z^2}$，即 $y^2$ 替换为 $y^2+z^2$，从而得到旋转曲面方程。
\end{note}
\bigskip

\begin{problem}
方程 $x^2+2z^2=4$ 在空间直角坐标系中表示\pickout{C}
\options
{椭球面}
{椭圆（平面曲线）}
{椭圆柱面，母线平行于 $y$ 轴}
{椭圆柱面，母线平行于 $z$ 轴}
\end{problem}
\begin{note}
柱面方程的特征是缺少某一个坐标变量。方程 $x^2+2z^2=4$ 中不含 $y$，说明该曲面由 xOz 平面上的椭圆沿 $y$ 轴方向平移生成，因此是母线平行于 $y$ 轴的椭圆柱面。
\end{note}
\bigskip

\end{document}
